\chapter{La dissonanza cognitiva: quando il cervello va in tilt}

\begin{flushright}
	\textit{``L'uomo è l'unico animale che arrossisce. O che ne ha bisogno.''}\\[6pt]
	\textsc{\index[main]{Twain, Mark}Mark Twain}
\end{flushright}

\bigskip
\bigskip

Vi è mai capitato di provare quella strana sensazione mentre scorrete lo schermo del telefono alle due di notte? Poche ore prima avevate promesso a voi stessi di dormire di più, magari dopo aver letto l'ennesimo articolo sui danni della privazione del sonno. Eppure eccovi lì, con la luce dello schermo a pochi centimetri dal viso, mentre fate scorrere foto di gatti e tramonti. Quella morsa allo stomaco, quel disagio sottile che non è esattamente senso di colpa ma gli somiglia molto, ha un nome preciso: dissonanza cognitiva.

Non è un concetto nuovo. Nel 1957 lo psicologo Leon Festinger lo descrisse per la prima volta in modo sistematico,\footnote{Festinger, L. (1957). \textit{A Theory of Cognitive Dissonance}. Stanford University Press, Stanford.} eppure gli esseri umani lo sperimentano da sempre. Probabilmente anche il primo uomo delle caverne che assaggiò una bacca velenosa, dopo aver avvertito la tribù di non farlo, conobbe quella sensazione. È il prezzo psicologico della contraddizione, il costo mentale dell'incoerenza, una sorta di mal di testa dell'anima che proviamo quando cerchiamo di far convivere pezzi che proprio non vogliono stare insieme. E come vedremo, dietro questa scoperta novecentesca si nasconde una delle più antiche ferite del pensiero umano.

\section{Un malinteso da chiarire: errore e disagio non sono la stessa cosa}

Prima di addentrarci nel labirinto, conviene fare una distinzione importante. I bias cognitivi, di cui abbiamo parlato nei capitoli precedenti, sono errori sistematici che il cervello commette quando elabora le informazioni. Sono automatici, inconsci, e li compiamo allegramente senza nemmeno accorgercene. Pensate al cervello come a un pilota automatico che prende scorciatoie: a volte funzionano benissimo, altre volte ci fanno sbattere contro un muro.

La dissonanza cognitiva è qualcosa di completamente diverso. Non è l'errore in sé, ma il disagio che proviamo quando ci accorgiamo che qualcosa non quadra. È come una sirena d'allarme che inizia a suonare quando il nostro comportamento tradisce i nostri valori, o quando due nostre convinzioni si scontrano come due treni lanciati sullo stesso binario in direzioni opposte. Ecco un modo per visualizzarla: il bias cognitivo è quando inciampate su un gradino, la dissonanza è il dolore al ginocchio che continua a ricordarvi di essere inciampati. Il bias è comprare l'ennesimo libro che non leggerete mai. La dissonanza è ciò che provate guardando la libreria stracolma di volumi impolverati mentre, nello stesso istante, cliccate ``acquista ora'' per il nuovo bestseller.

Questa distinzione, apparentemente tecnica, sfiora un nodo filosofico antichissimo. Perché il disagio della dissonanza presuppone qualcosa di tutt'altro che ovvio: che dentro di noi convivano voci diverse, capaci di giudicarsi a vicenda. Una parte fa, un'altra disapprova. Una agisce, un'altra arrossisce. È esattamente l'immagine con cui Platone, nella \textit{Repubblica}, aveva descritto l'anima: non un'unità compatta, ma un campo di battaglia, un'anima in guerra con se stessa, divisa tra una parte che desidera e una che sa.\footnote{Platone. \textit{La Repubblica}, libro IV, 439e e seguenti. Numerose edizioni.} La dissonanza, in fondo, è il rumore che fa quella guerra quando la sentiamo da dentro.

\section{Il fumatore moderno: una storia di contraddizioni}

Per capire davvero cos'è la dissonanza cognitiva non esiste esempio migliore del fumatore contemporaneo. Chiunque fumi nel ventunesimo secolo sa perfettamente che il fumo uccide. È scritto a lettere cubitali su ogni pacchetto, accompagnato da immagini che farebbero rabbrividire chiunque. Le prove scientifiche sono schiaccianti, documentate in migliaia di studi. Eppure continua ad accendere sigarette.

Dentro la sua testa convivono due affermazioni del tutto incompatibili: ``So che questa abitudine mi sta uccidendo lentamente'' e ``Continuo a farlo comunque''. È come tenere in mano due magneti con i poli dello stesso segno: si respingono, creano tensione, vogliono allontanarsi. Questa tensione è insopportabile per il cervello, che deve assolutamente trovare un modo per scioglierla. Ed è qui che la mente umana rivela una creatività straordinaria. Invece di affrontare la contraddizione nel modo più ovvio, cioè smettere, il cervello del fumatore si trasforma in un acrobata cognitivo, un maestro dell'arte della giustificazione.

C'è la minimizzazione del rischio: ``Mio nonno fumava due pacchetti al giorno ed è morto a novantacinque anni''. È l'esempio perfetto di come selezioniamo le prove che ci fanno comodo, ignorando la statistica e aggrappandoci all'unica eccezione come se fosse la regola. C'è l'aggiunta di giustificazioni positive: ``Il fumo mi rilassa'', ``Mi aiuta a concentrarmi'', ``È l'unico vizio che ho'', ``Tutti dobbiamo morire di qualcosa''. Ogni fumatore possiede la sua collezione personale di razionalizzazioni, costruite nel tempo come mattoni di un muro protettivo contro la verità. C'è la negazione vera e propria: ``Questi studi sono esagerati'', ``Il cancro è soprattutto una questione genetica'', ``Conosco gente che non ha mai fumato e si è ammalata lo stesso''. E poi, rarissima, c'è la modifica del comportamento, l'unica che affronti la dissonanza di petto: smettere davvero. Ma è difficile, dolorosa, richiede di ammettere a se stessi di aver sbagliato per anni. E così il cervello, nella stragrande maggioranza dei casi, sceglie la strada più comoda: cambiare le proprie credenze anziché il proprio comportamento.

Qui dobbiamo però fermarci su una sfumatura che cambia tutto. Siamo davvero sicuri che il fumatore non sappia? A guardarlo bene, sa benissimo. Recita anzi a memoria i rischi, li conosce meglio di molti non fumatori. Eppure si comporta come se non li conoscesse. Il filosofo francese Jean-Paul Sartre ha dato a questo strano stato un nome implacabile: \index[main]{mala fede}\textit{mala fede}.\footnote{Sartre, J. P. (1943). \textit{L'être et le néant}. Gallimard, Parigi.} La mala fede, per Sartre, è l'autoinganno attivo con cui ci nascondiamo una verità che in realtà possediamo. Non è ignoranza, perché per ingannarsi su qualcosa bisogna prima conoscerlo; e non è menzogna, perché il bugiardo sa di mentire mentre chi è in mala fede riesce quasi a crederci. È una posizione acrobatica dello spirito: sapere e insieme rifiutarsi di sapere. Il fumatore che enumera le sue giustificazioni non sta scoprendo nuove ragioni, sta costruendo attivamente il velo dietro cui smettere di vedere ciò che già vede. La razionalizzazione, da questo punto di vista, non è un guasto subito passivamente: è un'operazione che, in qualche angolo di noi, scegliamo di compiere.

\begin{figure}[htbp]
	\centering
	\captionof{figure}{La mente del fumatore è un campo di battaglia dove cognizioni incompatibili cercano disperatamente un equilibrio impossibile.}
	\label{fig:fumatore}
\end{figure}

\section{L'economia emotiva dell'armadio: quando lo sforzo crea valore}

Esiste un fenomeno curioso che gli psicologi hanno battezzato ``effetto IKEA''.\footnote{Norton, M. I., Mochon, D., e Ariely, D. (2012). The IKEA effect: When labor leads to love. \textit{Journal of Consumer Psychology}, 22(3).} In sostanza tendiamo ad attribuire più valore alle cose che abbiamo costruito con le nostre mani, anche quando il risultato è oggettivamente discutibile. Quella libreria assemblata in tre ore di sudore e imprecazioni vale oro ai vostri occhi, anche se è leggermente storta e vi sono avanzate tre viti di cui non avete mai capito la destinazione.

Perché accade questo? La risposta sta ancora una volta nella dissonanza. Quando investiamo tempo, energia e fatica in qualcosa, nella nostra mente si forma un'equazione implicita: ``Ho speso tre ore della mia vita per questa libreria, quindi deve essere fantastica''. L'alternativa, ammettere che quelle tre ore sono state uno spreco, è psicologicamente troppo costosa. Avete mai notato con quanto accanimento difendiamo le nostre scelte, anche quando sono palesemente sbagliate? Avete comprato un'auto che si rivela un disastro? Comincerete a magnificarne i pregi nascosti: ``Sì, consuma tanto, ma che linea''. Avete scelto un percorso di studi che si è rivelato noioso? All'improvviso diventa ``fondamentale per il mio futuro''. È il cervello che lavora a pieno regime per giustificare l'investimento già fatto e proteggervi dal dolore di ammettere l'errore.

Questo meccanismo ha implicazioni profonde. Spiega perché portiamo avanti relazioni che non funzionano (``Ci ho investito cinque anni, non posso mollare adesso''), perché restiamo in lavori che detestiamo (``Ho fatto troppi sacrifici per arrivare qui''), perché ci ostiniamo a finire libri noiosi (``Ho già letto duecento pagine, ormai devo''). Gli economisti lo chiamano fallacia dei costi irrecuperabili, ma alla radice c'è sempre lei, la dissonanza, che ci impedisce di tagliare i rami secchi della nostra vita.

\section{Il ciclo infinito: akrasia, dissonanza, razionalizzazione}

Ora che abbiamo capito cos'è la dissonanza, possiamo collocarla in un quadro più ampio: un ciclo che governa gran parte del comportamento umano, una danza in tre tempi che ripetiamo ogni giorno, spesso senza accorgercene.

Il primo tempo è l'\index[main]{akrasia}akrasia, la debolezza della volontà. Il termine viene dal greco antico e indica quella condizione in cui facciamo qualcosa pur sapendo che è sbagliato, che contraddice i nostri stessi obiettivi. Siete a dieta e passa un cameriere con una torta al cioccolato che sembra uscita da un dipinto rinascimentale. Il sistema limbico, la parte antica del cervello programmata da milioni di anni a cercare calorie e gratificazione immediata, entra in fibrillazione. La corteccia prefrontale, l'io razionale che conosce benissimo i vostri propositi, prova a resistere. Ma il sistema limbico è più veloce, più potente, evolutivamente più radicato. E così cedete: mangiate la torta. Non è un errore inconsapevole, è una scelta lucida di agire contro i propri interessi a lungo termine per una gratificazione immediata.

È proprio qui che vale la pena fermarsi, perché questo gesto così banale, una fetta di torta mangiata controvoglia, racchiude uno dei più antichi enigmi del pensiero occidentale. Fu Aristotele, nel settimo libro dell'\textit{Etica Nicomachea}, a porre la domanda nella sua forma più tagliente: com'è possibile che un uomo conosca il bene e scelga ugualmente il male?\footnote{Aristotele. \textit{Etica Nicomachea}, libro VII. Numerose edizioni.} Sembra una contraddizione logica, eppure accade di continuo. Il maestro di Aristotele, Socrate, aveva sostenuto che nessuno sbaglia volontariamente, che chi fa il male lo fa solo per ignoranza del bene. Aristotele, più attento all'esperienza concreta, non si accontentò di quella soluzione elegante e introdusse una distinzione preziosa. Esiste, sostenne, un sapere che possediamo soltanto in potenza, come chi dorme o è ubriaco: conosce le regole ma in quel momento non le ha davvero presenti. Nell'attimo della tentazione, la conoscenza del bene resta dentro di noi soltanto a parole, recitata come un attore recita un copione di cui non sente più il significato. Il desiderio immediato la mette a tacere. Tradotto nel linguaggio di oggi, diremmo che la corteccia prefrontale viene messa fuori gioco dall'urgenza limbica: ma l'intuizione di fondo, che la conoscenza intellettuale non basti a governare il comportamento, Aristotele l'aveva già colta ventitré secoli prima delle risonanze magnetiche.

Il secondo tempo è la dissonanza. Subito dopo, a volte nel giro di pochi secondi, scatta l'allarme. Nella vostra testa convivono ora due affermazioni inconciliabili: ``Sono una persona disciplinata che vuole dimagrire'' e ``Ho appena divorato cinquecento calorie di torta in meno di cinque minuti''. I due pensieri si scontrano, generano attrito, producono quel disagio che dà il nome al fenomeno. È come avere due file audio che suonano insieme su frequenze diverse: dissonante, alla lettera. La sensazione è profondamente sgradevole, e il cervello vuole farla sparire al più presto.

Ed ecco il terzo tempo, la razionalizzazione. Il cervello cerca disperatamente un modo per cancellare il disagio, e la via più semplice, quella che non chiede di ammettere l'errore né di cambiare comportamento, è distorcere la realtà. Così vi dite: ``In fondo oggi ho fatto le scale invece dell'ascensore'', ``Una fetta ogni tanto non cambia nulla'', ``Domani recupero con una camminata''. Queste giustificazioni non devono essere logiche né coerenti: devono solo funzionare abbastanza da ridurre la dissonanza, permettendovi di conservare insieme il comportamento sbagliato e l'immagine positiva di voi stessi. È un gioco di prestigio cognitivo: far sparire la contraddizione senza risolvere il problema. E il ciclo ricomincia. Domani, dopodomani, la settimana prossima. Akrasia, dissonanza, razionalizzazione, ancora e ancora.

\begin{figure}[htbp]
	\centering
	\includegraphics[width=0.9\textwidth]{images/disso.png}
	\captionof{figure}{Il ciclo di akrasia, dissonanza e razionalizzazione illustrato nelle sue tre fasi ricorrenti.}
	\label{fig:disso}
\end{figure}

\section{Le radici profonde: perché l'evoluzione ci ha messi in difficoltà}

A questo punto potreste domandarvi: ma perché il cervello funziona in un modo così controproducente? Se siamo tanto bravi a ingannare noi stessi, come abbiamo fatto a sopravvivere come specie? La risposta sta nel nostro passato, in quei milioni di anni in cui i nostri antenati erano cacciatori e raccoglitori. In quel contesto, i meccanismi che oggi ci creano problemi erano preziosi vantaggi.

L'akrasia, il privilegiare il presente, aveva perfettamente senso. In un mondo di scarsità, dove il cibo era incerto e i predatori una minaccia costante, mangiare tutto ciò che si trovava oggi era la strategia vincente: chi rinunciava per ``risparmiare per il futuro'' rischiava di morire di fame, o di finire tra le fauci di un leone, prima ancora di vedere quel futuro. Il nostro cervello è ancora cablato per quella realtà. La dissonanza, a sua volta, si è sviluppata come meccanismo per mantenere coerenza nelle credenze, cosa fondamentale per agire in fretta in un ambiente pericoloso: non potevi permetterti di restare paralizzato dai dubbi, ti serviva un sistema di convinzioni stabile, anche se imperfetto. E la razionalizzazione protegge l'autostima, che non è un lusso psicologico ma un requisito di sopravvivenza sociale: negli ambienti ancestrali essere esclusi dal gruppo significava morte quasi certa, e mantenere un'immagine di sé sufficientemente positiva da restare integrati era vitale.

Il problema è che oggi viviamo in un mondo radicalmente diverso da quello per cui questi meccanismi furono progettati. Abbiamo cibo spazzatura disegnato per forzare i nostri circuiti della ricompensa, piattaforme costruite scientificamente per generare dipendenza, carte di credito che rendono invisibili le conseguenze delle nostre spese, servizi che con il loro ``prossimo episodio tra cinque secondi'' sfruttano la nostra incapacità di resistere alla gratificazione immediata. E così il ciclo si ripete: agiamo male, ci sentiamo in colpa, ci giustifichiamo, e ricominciamo, intrappolati in schemi che sappiamo dannosi ma che fatichiamo enormemente a spezzare, come criceti in una ruota che gira all'infinito.

C'è però un'ombra ancora più inquietante in tutto questo, che riguarda non la sopravvivenza ma la morale. Nel pieno dell'Illuminismo, Immanuel Kant aveva immaginato un essere umano capace di agire per puro dovere, seguendo l'imperativo categorico, ossia la regola di comportarsi sempre secondo una massima che potremmo desiderare valesse come legge universale, a prescindere da ogni inclinazione.\footnote{Kant, I. (1788). \textit{Kritik der praktischen Vernunft}. Riga.} Era una scommessa altissima sulla ragione: la volontà buona, sosteneva Kant, deve poter zittire i desideri con la sola forza del principio. Il ciclo akrasia, dissonanza, razionalizzazione è la cronaca minuziosa di quanto raramente ci riusciamo. Non perché l'ideale kantiano sia sbagliato, ma perché l'animale che dovrebbe realizzarlo è lo stesso che, davanti alla torta, sente la corteccia prefrontale spegnersi. Kant descriveva chi vorremmo essere; la neuroscienza descrive chi siamo. E lo spazio fra le due descrizioni è esattamente il territorio della dissonanza.

\section{Consapevolezza: necessaria ma non sufficiente}

Riconoscere questo meccanismo è importante, fondamentale persino. È il primo passo per interrompere il ciclo. Ma sarebbe ingenuo credere che la sola consapevolezza basti. Lo aveva intuito Spinoza, quando scrisse che la vera conoscenza del bene e del male non riesce a frenare un sentimento in quanto è vera, ma solo in quanto è essa stessa vissuta come un sentimento.\footnote{Spinoza, B. (1677). \textit{Ethica}, parte IV, proposizione XIV. Pubblicazione postuma, Amsterdam.} In parole più semplici: sapere che qualcosa fa male non basta per smettere di farlo. Milioni di fumatori ne sono la prova vivente. La conoscenza è un punto di partenza, non un punto di arrivo.

Ed è sorprendente constatare quanti grandi pensatori abbiano battuto, da angolazioni diverse, su questa stessa porta. Aristotele si chiedeva come l'uomo sapiente potesse comportarsi da stolto. Kant cercava di costruire una ragione capace di comandare ai desideri. Sartre denunciava la fuga in mala fede da ciò che già sappiamo. Freud, qualche decennio dopo Sartre ma prima delle neuroscienze, avrebbe parlato di un Io perennemente assediato dalle pulsioni dell'Es. Tutti, con linguaggi lontani, hanno cercato di rispondere alla stessa domanda: perché l'uomo che sa il bene fa il male? Nessuno l'ha risolta del tutto. Ma ognuno ne ha illuminato un lato diverso, e oggi le neuroscienze aggiungono il loro tassello senza pretendere di chiudere il discorso. La conoscenza dei nostri trucchi resta preziosa: capire come funziona il cervello, riconoscere i momenti in cui stiamo razionalizzando invece di ragionare, tutto questo ha un valore reale. Ma poi serve altro, perché contro un avversario così antico le buone intenzioni non bastano: serve modificare l'ambiente, costruire sistemi di supporto, sviluppare abitudini alternative, a volte cercare aiuto professionale.

Perché alla fine, come abbiamo visto, non stiamo combattendo contro un piccolo difetto del nostro modo di pensare. Stiamo combattendo contro milioni di anni di evoluzione, contro meccanismi profondamente radicati che un tempo ci hanno salvato la vita e che oggi, in un mondo del tutto diverso, rischiano di rovinarla. La buona notizia è che siamo anche l'unica specie capace di accorgersene: di studiare il proprio cervello, di riconoscere i propri limiti e, talvolta, di inventare modi creativi per aggirarli. Non è una battaglia facile, ma almeno adesso sappiamo contro cosa stiamo combattendo. E questa, forse, è già una piccola vittoria.